\section{Safety and Science}
\sectionkicker{SAFETY}
\begin{wideflow}
{\Large\bfseries 守りの制度、探究の公開、模擬の測定――三件の境界線\par}
\smallskip
{\color{SurveyMuted}サイバー計画の発足、数学成果の公開、模擬評価の三件を、効果・入手・実環境の未確定とともに整理する。\par}
\end{wideflow}

\subsection*{守りと探究の三件――実測なき期待と未公開の成果}

Anthropicは10月8日、Cyber Missionを開始した。柱は二つで、11社の創設パートナーを擁するCIDPと、OSS Scannerという無償のオプトイン検査器だ。Scannerはモデル生成の概念実証レポートを出し、人間の査読を挟まず、真陽性率90\%超を見込むとされる。制度の発足としては確定情報であり、参加条件と提供形態の層では具体的に語れる。\cite{cybermission}

ただし期待と実測は別物だ。90\%超は見込みの数字であり、測定された成果ではない。Scanner本体の詳細投稿もCIDPの申込フォーム内部も未取得で、実際にどれだけ守れるかは未知数だ。制度の存在を効果の実証と読み替えないことが、この記事の第一の留保になる。\cite{cybermission}

\subsection*{数学の成果公開――モデルは出さず、中身を出す}

OpenAIは10月6日、社内の最前線モデル（名称非公開）に由来する広範な数学成果を公開した。GitHubのopenai／mathリポジトリ、Leanによる部分形式化、10件の推論要約、Pro利用での計算統計（平均約3時間）が列挙されている。成果物の公開としては確定情報であり、形式化の範囲と要約の存在は開発者・研究者の参照点になる。\cite{mathrelease}

ここで区別すべきは、成果の公開とモデルの公開だ。モデル本体は未公開で、責任ある公開に向けた作業が継続中とされ、期日はない。リポジトリの内容物や論文数、Leanファイルの数え上げも未列挙で、ライセンスもページにない。数学の前進として読むべきもので、モデル入手の見通しとして読むべきものではない。\cite{mathrelease}

\subsection*{模擬で測る――Ironcladのコンピュータ利用評価}

OpenAIとIroncladの10月6日の研究は、11件の契約タスクによるコンピュータ利用の評価だ。条件はEDGAR模擬で、顧客・社内・非公開データは使われていない。結果は成功率でAstra Max 55.0\%に対しSol High 41.6\%、差は13.4ポイント、相対改善率は約32\%とされる。所要時間は37.0分から19.2分へほぼ半減（約48\%減）とされる。いずれも模擬条件でのベンダー測定であり、暫定的な研究段階の数字だ。絶対差と相対差の書き分けに注意したい。13.4ポイントは二つの成功率の差、約32\%は低い方を基準にした相対の伸びであり、別の量だ。\cite{ironclad}

模擬の意味を正確に押さえたい。EDGAR模擬の契約タスクという母集団は既知だが、採点ルーブリックやタスク別の得点、強化学習の詳細はページにない。実環境への導入効果として読むのは誤りで、あくまで模擬条件での相対比較だ。Ironcladの製品発表は本ページ外の別事象であり、混同してはならない。\cite{ironclad}

三件を通覧すると、共通するのは「確定した制度・公開・模擬」と「未確定の効果・入手・実環境」という二層構造だ。Cyberは制度確定・効果未測定、数学は成果確定・モデル未公開、Ironcladは模擬確定・実環境未測定という対応になる。この対応表自体が本記事の読みどころであり、三件を一つの総合性能ランキングにまとめることはしない。模擬評価から判断できるのは、模擬条件の範囲での相対比較までであり、導入効果の予測ではない。\cite{cybermission}\cite{mathrelease}\cite{ironclad}

未取得の三点をまとめて示す。Scanner詳細と申込内部、リポジトリの数え上げとライセンス、ルーブリックとタスク別得点である。いずれも次の取材・取得の条件であり、推測での穴埋めはしない。囲いの中の事実と囲いの外の未知を分けて示すことが、安全と科学の主題に対する正確さだ。\cite{cybermission}\cite{mathrelease}\cite{ironclad}

実務への含意は控えめに述べる。セキュリティ担当者はScannerの無償オプトインを試用候補に載せられる段階にあり、数学・自動化の関心層はリポジトリと模擬条件を参照点にできる段階にある。いずれも導入判断ではなく、次の検証への入口だ。\cite{cybermission}\cite{mathrelease}\cite{ironclad}

\begin{technicalnote}[Note]
期待性能と実測、成果公開とモデル公開、模擬と実環境をそれぞれ区別した。成功率の差は13.4ポイント、相対改善は約32\%。いずれも独立再現なし。
\end{technicalnote}
